% =====================================================================
\chapter{Conclusion and Future Work}\label{ch:conclusion}
% =====================================================================

% \begin{plan}[Chapter plan --- Conclusion]
% \pages{2--3 pages.} No new results, no new citations except in future work.
% \end{plan}

\section{Conclusion}\label{sec:conc-summary}

This thesis developed and evaluated a protocol for auditing gender-related
responses in vision language models through their next-token probability
distributions. The evaluation covered five models from two families, ranging
from 2B to 90B parameters, across three experimental phases and 25,450
measurements. The experiments examined responses to inventory statements,
changes associated with face images and downstream judgements of image search
results. Human ratings provided a comparison for the downstream task,
although differences in presentation limit the interpretation of that comparison.

The protocol obtains each measurement from a single forward pass, pooling the
probability assigned to alternative surface forms of the intended responses.
Its implementation specifies the chat template, assistant prefill and position
at which the distribution is read. These choices proved necessary for
obtaining measurements of the intended response. Moving the rating cue into
the assistant turn increased captured mass over the seven digits from
$0.0043$ to $0.994$, demonstrating how strongly the measurement can depend on
input construction.

Captured mass provided a practical check on whether the selected response set
adequately represented the model's output distribution. Six measurement faults
were identified during the project, none of which produced an execution error.
In one case, reading only lowercase response tokens captured $0.09\%$ of the
distribution while still producing a plausible aggregate score. In another,
a hardware fault that zeroed hidden states produced an expected rating of
exactly $4.000$. These cases show why apparently reasonable numerical outputs
require checks of both probability coverage and output diversity.

The downstream experiments identified a substantial difference in scale
reversal robustness between the visual and text-only conditions for the two
largest models. Without an image, Llama-90B and Qwen-72B each reproduced their
normalised judgement at a robustness rate of $0.900$. With an image, these
rates fell to $0.044$ and $0.156$, giving $\dRR=-0.856$ and $-0.744$,
respectively, with bootstrap intervals excluding zero. The smaller models had
low robustness in both conditions, so their smaller differences did not
indicate more reliable judgements. The shared pattern in the two largest
models warrants further investigation, although the evaluated conditions do
not isolate its internal cause.

The binary inventory probes also showed that the effect of visual input
differed between the evaluated model families. Both Llama models assigned
substantially less probability to the yes and no response set when an image
was present, including a featureless grey patch. The Qwen models maintained
high coverage across the corresponding conditions, with the contrast extending
to the face image experiments. For example, captured mass was $0.140$ for
Llama-11B and $0.998$ for Qwen-7B when real faces were shown. This behaviour
is consistent with model-dependent defensiveness, but probability outside the
intended response set does not establish refusal without examining the
alternative outputs.

The analysis further demonstrated that the conventional aggregate inventory
score is sensitive to a model's general tendency to agree. Because the
inventory contains sixteen positively keyed items and six reverse-keyed items,
uniform assent contributes positively to the aggregate even without meaningful
discrimination between items. The decomposition showed that Llama-90B's score
of $+0.153$ was close to the $+0.160$ predicted by its mean assent rate alone.
Polarity separation addresses this limitation by measuring the difference in
mean assent between the two item groups, remaining invariant to uniform shifts
in agreement. The resulting comparison describes item discrimination and
should not be interpreted directly as a ranking of models by sexism.

The face image experiments provided a separate finding concerning the
attribution of benevolent traits. Four models attributed these traits more
readily to photographs of women, while the smallest model showed the opposite
direction. For Qwen-72B, the positive gap between female and male faces reached
$+0.054$, with $p<0.001$.
These differences demonstrate the importance of examining the subscales
separately, since aggregation can obscure patterns that differ across item groups.

Taken together, the findings support three practical requirements for
probability-based model audits. Captured mass should accompany the reported
measurements so that their coverage remains visible. Aggregate scores should
distinguish item-level discrimination from a general tendency to agree.
Rating protocols should also be tested for robustness to equivalent scale
definitions before their outputs are interpreted as stable judgements.
These checks strengthen the interpretation of the measurements while
preserving the distinction between observed model responses and claims about
underlying attitudes.

% \begin{plan}
% One paragraph per contribution of \cref{sec:intro-contrib}, each with its
% single strongest number. End with one sentence on what an auditor should do
% differently because of this work (report coverage; use style-invariant scores;
% test robustness to the stated scale).
% \end{plan}

\section{Future work}\label{sec:conc-future}

A useful extension would separate responses to gendered referents from
responses to the inventory statements themselves. The present design could
be expanded by asking matched items about a woman, a man and a person,
introducing the referent as an additional experimental factor. This comparison
would help determine whether a response concerns women specifically or
reflects a broader attribution to people, without requiring a photograph to
establish the referent (\cref{sec:disc-threats}).

The downstream task could also be brought into closer correspondence with
the human evaluation. Presenting all nine search results as a single composite
would allow the model to judge the same grid that the human participants
evaluated. This would reduce the difference between individual-image model
ratings and grid-level human judgements, making deviations from the human
baseline easier to interpret. Additional non-person queries with comparable
human baseline ratings would also strengthen the control comparison beyond
the single balloon query used here.

The human comparison could be extended by calculating the position index
$\lambda$ from the released participant responses (\cref{eq:lambda}). This
index would locate each model relative to both the low and high ASI groups,
providing information that distance from the low ASI baseline alone cannot
supply. Such an analysis would still require caution where model ratings
remain sensitive to scale orientation, since the comparison assumes that the
reported judgements are sufficiently stable.

Further work could investigate which model components contribute to the
redistribution of probability when visual input is introduced. Activation
analysis and controlled interventions could examine whether the observed
changes arise during visual processing, multimodal integration or subsequent
language processing. The approach used by Arditi et al.\ \cite{arditi2024refusal}
to identify a direction associated with refusal provides one possible
methodological starting point. However, the present observations do not
establish that the probability shifts share the same mechanism.

A contemporary collection of results for the same ten queries would permit an
evaluation of changes since the original 2018 dataset. Comparing historical
and current results under a consistent protocol would help distinguish changes
in search content from differences in model judgement. The comparison would
require recording the collection conditions carefully, since search results
can vary with location, personalisation and collection time.

Finally, a forced-choice version of the inventory could test whether the
observed item patterns persist under a different response format. Asking models
to choose between matched alternatives would reduce reliance on direct assent
and provide a complementary test of the response style effects identified here.
The alternatives and their ordering would require careful control, since
changing the format can introduce preferences of its own. Agreement across
these formats would provide stronger evidence that a finding reflects
discrimination between item meanings rather than sensitivity to a particular
prompt construction.

% \begin{plan}
% \begin{itemize}
%   \item \textbf{Gendered referents} (woman / man / person) as a fourth factor
%         in Phase~1 --- addresses the main construct threat.
%   \item \textbf{Grid-level stimuli} in Phase~2 to match the human task.
%   \item \textbf{High-ASI baselines and the position index} $\lambda$
%         (\cref{eq:lambda}).
%   \item \textbf{Contemporary 2026 stream} for temporal drift (proposal).
%   \item \textbf{Attribution}: cross-attention or activation analysis of what
%         moves mass away from \{yes, no\} (cf.\ \cite{arditi2024refusal}).
%   \item \textbf{Forced-choice probes} that remove the yes-rate by
%         construction.
% \end{itemize}
% \end{plan}
